% Section 2.8, translated from sv/chapters/02/summary.tex.

\section{Summary}
\label{c2:sec:sammanfattning}

\begin{booktable}{@{}p{12.5em}L@{}}
  \thead{Rule} & \thead{Formula} \\
  \midrule
  Distributive law & $a(b + c) = ab + ac$ \\
  Square of a sum & $(a + b)^2 = a^2 + 2ab + b^2$ \\
  Square of a difference & $(a - b)^2 = a^2 - 2ab + b^2$ \\
  Difference of two squares & $(a + b)(a - b) = a^2 - b^2$ \\
  Common factor & $ab + ac = a(b + c)$ \\
  Power & $P = UI = RI^2 = \frac{U^2}{R}$ \\
\end{booktable}

\needspace{6\baselineskip}
\subsection*{What you should be able to do}
After this chapter, you should be able to:
\begin{itemize}
  \item Understand what an algebraic expression is and how variables are used.
  \item Simplify algebraic expressions.
  \item Expand and factorise expressions.
  \item Use the rules to simplify and speed up calculations in electrical circuits.
\end{itemize}

\testadigsjalv
\begin{enumerate}
  \item Factorise $9R^2 - 4$.
  \item Simplify $2(3I + 1) - (I - 5)$.
  \item A heating element with $R = 20\,\Omega$ is supplied first with $U_1 = 52\,\text{V}$ and
    then with $U_2 = 48\,\text{V}$. Use the difference of two squares to calculate the difference
    in power
    \[
      P_1 - P_2 = \frac{U_1^2 - U_2^2}{R}
    \]
    without a calculator.
\end{enumerate}

\needspace{6\baselineskip}
\subsection*{Next chapter}
\Kapref{3} is about linear equations and inequalities: how they are solved, how the solution of an
inequality is drawn on the number line and how equations with absolute values are handled.
